\section{Discussion, Future Work, and Conclusion}
\label{sec:discussion}

\subsection{Synthesis of Results}
\label{subsec:disc-synthesis}

Relational Superposition Theory offers a coherent mathematical vocabulary for
concurrent non-exclusive dating arrangements without forcing them into
ill-fitting total orders. Starting from the observation that multi-role
functional specialization destroys scalar Pareto comparability (\cref{prop:order-failure}),
we represented partner portfolios as normalized state vectors over a role basis
and established three interlocking mechanisms:
\begin{enumerate}[nosep]
  \item \textbf{Static indeterminacy (\cref{prop:nobody-first,thm:rank-role-inequality}).}
        Prior to collapse, no partner holds unit primary weight, and rank and
        role dispersions obey $\Delta R \cdot \Delta L \ge \kappaS$ whenever the
        rank and role operators fail to commute.
  \item \textbf{Dynamical maintenance and sudden failure (\cref{prop:stationary,prin:poke-rule,thm:percolation-threshold}).}
        Routine weekday, weekend, and stress cycles maintain positive rank
        entropy indefinitely in stationarity, whereas a poke forces irreversible
        projective collapse (\cref{prop:irreversibility}) or a network-wide
        percolation cascade above critical friend-graph density $q_c$.
  \item \textbf{Structural drivers (\cref{thm:optionality-monotonicity,prop:symmetry}).}
        The prevalence of superposition scales monotonically with viable option
        count $n(A)$ rather than individual disposition or gender, placing
        both the baseline configuration and its mirror case on identical formal
        footing.
\end{enumerate}

\subsection{Relation to Preference and Matching Literature}
\label{subsec:disc-literature}

In classical two-sided matching~\citep{gale1962college,roth1990two}, instability
arises when blocking pairs prefer one another over assigned partners under
fixed total orders. In situationships, instability is internal to a single
agent's preference structure: the focal agent avoids committing to a total order
precisely because doing so forfeits specialized complementary roles held by
non-primary partners. Rather than blocking-pair instability, the primary
failure mode is informational percolation across the shared social graph,
linking individual preference indeterminacy directly to network threshold
phenomena~\citep{granovetter1978threshold,watts2002simple}.

\subsection{Directions for Future Work}
\label{subsec:disc-future}

Four concrete extensions merit further investigation:
\begin{enumerate}[nosep]
  \item \textbf{Empirical estimation of $\kappaS$ and $p$.} Deploying the
        non-intrusive observational cadence protocol in \cref{app:instrument}
        on anonymized messaging logs would test the option-mediation prediction
        of \cref{conj:optionality} and yield empirical ranges for $\kappaS$.
  \item \textbf{Endogenous poke feedback.} Coupling the Markov attention-drift
        model of \cref{sec:dynamics} with neglect-triggered poke hazards
        $\lambda_i(w_i(t))$ would unify pre-poke drift and collapse timing into
        a single continuous-time jump-diffusion process.
  \item \textbf{Multi-partner simultaneous pokes.} Extending the projective
        collapse operator of \cref{sec:poke} to simultaneous non-commuting
        queries from two partners ($G_i$ and $G_j$ poking within the same
        silence window $\tau_s$) introduces order-dependent collapse branching.
  \item \textbf{Continuous role manifolds.} Replacing the discrete $m$-role
        basis $\roleset$ with a continuous latent functional space learned from
        interaction embeddings would remove basis-granularity sensitivity.
\end{enumerate}

\subsection{Conclusion}
\label{subsec:disc-conclusion}

When a person maintains multiple concurrent partners across distinct functional
roles, asking ``who is first?'' presupposes a classical ordering that the
system has not constructed. Relational Superposition Theory replaces that
presupposition with an explicit state-space model in which rank indeterminacy,
role specialization, non-commutation, poke-induced collapse, social-graph
decoherence, and optionality scaling follow from unified mathematical axioms.
Until a poke collapses the state, no partner is first or second, and rank and
role cannot be resolved together.
